\section*{Hoofdstuk \ref{ch:g6:exploring-our-environment} --- Onze omgeving verkennen}

\begin{solution}{exo:g6:exploring-our-environment:1}
\omterm{def:g6:exploring-our-environment:components}{Levende wezens} (mos, pissebedden); \omterm{def:g6:exploring-our-environment:components}{minerale materie} (de steen van de
muur, de plas); sporen en resten van \omterm{def:g6:exploring-our-environment:components}{levende wezens} (dode \omterm{def:g1:plants-around-us:parts}{bladeren}, het
slakkenhuis).
\end{solution}

\begin{solution}{exo:g6:exploring-our-environment:2}
\omterm{def:g6:exploring-our-environment:conditions}{Temperatuur}, met een thermometer (in graden Celsius); de \omterm{def:g6:exploring-our-environment:conditions}{vochtigheid} van
lucht en bodem; \omterm{def:g6:exploring-our-environment:conditions}{licht}, van volle zon tot diepe schaduw.
\end{solution}

\begin{solution}{exo:g6:exploring-our-environment:3}
Levend: de varen, de duizendpoot. Mineraal: de plas, de steen van de
muur. Sporen: het slakkenhuis, de vogelpoep.
\end{solution}

\begin{solution}{exo:g6:exploring-our-environment:4}
Zestien graden (\qty{38}{\degreeCelsius} tegen
\qty{22}{\degreeCelsius}) en volle zon tegen diepe schaduw (en ook droog
tegen vochtig). Het onkruidje dat droogte en felle zon verdraagt past bij
de zuidvoet; mos, \omterm{prop:g4:plants-without-seeds:damp}{varens} en pissebedden passen bij de koele vochtige hoek.
\end{solution}

\begin{solution}{exo:g6:exploring-our-environment:5}
Zijn bedekking laat water ontsnappen, dus verliest zijn lichaam in droge
zon dodelijk snel zijn water. Daarom vind je pissebedden in vochtige,
donkere \omterm{def:g6:exploring-our-environment:microhabitat}{microhabitats} --- onder stenen, onder dood hout.
\end{solution}

\begin{solution}{exo:g6:exploring-our-environment:6}
Bijvoorbeeld: een aangeknaagde hazelnoot --- de bosmuis; een spinnenweb
--- de spin; sporen in de modder --- de reiger (uitwerpselen, \omterm{def:g1:animals-around-us:covering}{veren} en
\omterm{def:g1:animals-around-us:covering}{schelpen} doen hetzelfde).
\end{solution}

\begin{solution}{exo:g6:exploring-our-environment:7}
Een gesloten doos, half vochtig papier en half droog, met de pissebedden
losgelaten langs de middenlijn; een paar minuten later zit bijna alles op
de vochtige kant (net zo bij schaduw tegen \omterm{def:g6:exploring-our-environment:conditions}{licht}). Het laat zien dat elk
\omterm{def:g1:animals-around-us:animal}{dier} neerstrijkt waar de \omterm{def:g6:exploring-our-environment:conditions}{omstandigheden} bij zijn lichaam passen --- de
kaart van de groep tekent zichzelf, zonder opdrijven.
\end{solution}

\begin{solution}{exo:g6:exploring-our-environment:8}
Omdat alleen metingen die op dezelfde manier gedaan zijn vergeleken
kunnen worden: een verschil moet van de locaties komen, niet van het
meten. Het is de veldversie van de eerlijke proef --- één ding varieert
(de locatie), al het andere blijft gelijk.
\end{solution}

\begin{solution}{exo:g6:exploring-our-environment:9}
Bovenkant van de steen: verlicht, droog, warm --- korstmos (met de
zonnende vlieg als bezoeker). Onderkant van de steen: donker, vochtig,
koel --- pissebedden (of de duizendpoot). Vijf centimeter uit elkaar,
twee bijpassende verzamelingen bewoners.
\end{solution}

\begin{solution}{exo:g6:exploring-our-environment:10}
Stap 1: markeer een duidelijke vierkante meter. Stap 2: meet --- de
thermometerstand opgeschreven, de \omterm{def:g6:exploring-our-environment:conditions}{vochtigheid} van de bodem gevoeld en
genoteerd, de lichtklasse genoteerd. Stap 3: zoek goed, stenen opgetild
en teruggelegd, alle drie de \omterm{def:g5:biodiversity:species}{soorten} bestanddelen opgesomd. Stap 4:
noteer ter plekke --- aantallen, niet ``wat''. Stap 5: met één locatie
was er niet eens een vergelijking mogelijk; onderzoek een tweede locatie
op dezelfde manier.
\end{solution}

\begin{solution}{exo:g6:exploring-our-environment:11}
Verklaring: klimop gedijt binnen het bereik van \omterm{def:g6:exploring-our-environment:conditions}{omstandigheden} van de
noordmuur (schaduw, vocht) en niet binnen dat van de zuidmuur (felle zon,
droogte). Toets: meet beide muurvoeten --- \omterm{def:g6:exploring-our-environment:conditions}{temperatuur}, \omterm{def:g6:exploring-our-environment:conditions}{vochtigheid},
\omterm{def:g6:exploring-our-environment:conditions}{licht} --- en vergelijk met de bekende voorkeuren van klimop; of zoek
klimop op andere muren en noteer welke oriëntatie die muren delen.
\end{solution}

\begin{solution}{exo:g6:exploring-our-environment:12}
Haal het verschil in wanden weg: doe de keuzeproef in een ronde doos
(overal dezelfde wand), waarbij je de vochtige en droge helft tussen de
rondes verwisselt --- vochtig in het oosten in ronde 1, vochtig in het
westen in ronde 2. Als de pissebedden het vocht volgen waar je het ook
neerlegt, vallen de wanden af; alleen de onderzochte omstandigheid
veranderde tussen de rondes, al het andere bleef gelijk.
\end{solution}

\begin{solution}{pb:g6:exploring-our-environment:1}
\textbf{1.} De hoge heg: die geeft sloot N schaduw (en beschutting),
terwijl sloot Z de hele dag op de zon ligt.

\textbf{2.} \omterm{def:g6:exploring-our-environment:conditions}{Temperatuur}: \qty{19}{\degreeCelsius} tegen
\qty{31}{\degreeCelsius}. \omterm{def:g6:exploring-our-environment:conditions}{Vochtigheid}: de grond van sloot N is vochtig
(plakt), die van sloot Z droog (stof). \omterm{def:g6:exploring-our-environment:conditions}{Licht} maakt het drietal compleet:
schaduw tegen volle zon.

\textbf{3.} De \omterm{def:g6:exploring-our-environment:conditions}{omstandigheden} veranderen door de dag heen: om zes uur
's avonds is elke locatie koeler en is het \omterm{def:g6:exploring-our-environment:conditions}{licht} zwakker, dus zou het
gemeten verschil locatie met uur vermengen. Hetzelfde uur houdt de
vergelijking eerlijk.

\textbf{4.} Zodat de zoekinspanning gelijk is: meer oppervlak of meer
tijd op de ene locatie zou daar meer bewoners opleveren --- een oneerlijke
toets van welke sloot rijker is.

\textbf{5.} Beide lijsten noemen alleen \omterm{def:g6:exploring-our-environment:components}{levende wezens}. Ontbreken: de
minerale bestanddelen (bodem, stenen, water) en de sporen --- die een
behoorlijke inventarisatie ook moet vastleggen.

\textbf{6.} Pissebedden: sloot N --- ze hebben vocht en schaduw nodig,
anders drogen ze uit. De pad: sloot N --- een vochtige \omterm{def:g1:the-five-senses:sense}{huid}, die vochtige
schuilplaatsen vraagt. De hagedis: sloot Z --- ze warmt haar lichaam in
de zon en jaagt waar ze kan zonnen.

\textbf{7.} Of er leven daar slakken die zich verstoppen (misschien
alleen actief in vochtige nachturen), of het huisje is een rest van
elders of uit een ander seizoen --- een spoor bewijst aanwezigheid
\emph{op enig moment}, geen huidige bewoning.

\textbf{8.} Donker, vochtig, koel. De regel van
\cref{prop:g6:exploring-our-environment:decide}: elke plek wordt bewoond
door de \omterm{def:g5:biodiversity:species}{soorten} bij wier bereik de \omterm{def:g6:exploring-our-environment:conditions}{omstandigheden} passen.

\textbf{9.} N heeft mos en padden omdat de schaduw het er koel en
vochtig houdt, wat hun lichaam nodig heeft; Z heeft hagedissen en
\omterm{def:g1:plants-around-us:plant}{planten} met wasachtige \omterm{def:g1:plants-around-us:parts}{bladeren} omdat de volle zon het er heet en droog
maakt, wat een zonnende jager en een droogtebestendige \omterm{def:g1:plants-around-us:plant}{plant} past.

\textbf{10.} Zonder de heg krijgt sloot N volle zon: heter, droger ---
\omterm{def:g6:exploring-our-environment:conditions}{omstandigheden} die naar die van de huidige sloot Z opschuiven. Verwacht
dat de vochtliefhebbers (mos, \omterm{prop:g4:plants-without-seeds:damp}{varens}, pissebedden, de pad) teruglopen en
dat zonliefhebbers (droge grassen, sprinkhanen, misschien de hagedis)
hun intrek nemen.

\textbf{11.} Omdat het tweede onderzoek in één ding van het eerste moet
verschillen --- de afwezigheid van de heg. Nieuwe locaties, andere
instrumenten of een ander seizoen zouden eigen verschillen toevoegen, en
dan zou de verandering niet meer aan de heg toegeschreven kunnen worden:
de regel van de eerlijke proef op de schaal van een landschap.

\textbf{12.} Elke pissebed loopt rond en keert tot zijn lichaam zich
prettig voelt, dus komt de groep uit waar het vocht zit --- de posities
van de \omterm{def:g1:animals-around-us:animal}{dieren} leggen de \omterm{def:g6:exploring-our-environment:conditions}{omstandigheden} vast, net zo betrouwbaar als
thermometerstanden. Niemand drijft ze op; het eigen comfort van elk
individu tekent de kaart, en juist daarom kunnen bewoners als metingen
gelezen worden.
\end{solution}
